% Disjoint FV authoring packet; source-version and whole-source duties remain separate.
\chapter{Feferman--Vaught Composition and Boolean Powers}
\label{sec:p03-feferman-vaught-boolean-powers}

The truth of a formula in a direct product is not, in general, the assertion
that the same formula is true in every factor. The Feferman--Vaught method
instead produces finitely many factor formulas and a first-order condition
on their truth sets in the Boolean algebra of indices. The existential step
must construct compatible local witnesses, not merely name a composition
theorem. The partition/patching mechanism is explicitly treated in
\cite[\S2.5, Lemma~2.9 and Theorem~2.10]{DaquinoMacintyre2023FVPreprint};
that source's ring/idempotent presentation is distinct from the arbitrary
signature construction given below. The historical general source is
\textcite{FefermanVaught1959}, and the bounded Boolean-power presentation is
\cite[IV\S5]{BurrisSankappanavar2012}.

\section{Products, local truth sets, and the Boolean language}

\begin{hbconvention}{def:p03-fv-standing}{Domains, sorts, and Choice}
Let \(L\) be a set-sized first-order signature with finite-arity functions
and relations and logical equality. Many sorts are allowed, with the same
sort and arity declarations in every factor. Every sort of every
\(L\)-structure \(M_i\), \(i\in I\), is nonempty, and \(I\) is a set.
The full-product theorem is proved in \(\mathsf{ZF}+\mathsf{AC}\).
Choice supplies nonempty product sorts and the coordinate witness functions
in the existential step; it is a metatheoretic assumption, not an axiom in
the factor theories. Only the finitely many symbols occurring in the input
formula affect its reduction; neither the signature nor the index set is
required to be finite.

Empty products are included. Their sort domains are the singleton empty
tuples. Accordingly the Boolean algebra of indices may be trivial:
\(\mathcal P(\varnothing)\) has \(0=1\). Boolean logical falsity below
is written \(0\ne0\), not \(0\ne1\), so that this edge case does not
change the meaning of falsity.
\end{hbconvention}

\begin{hbdefinition}{def:p03-fv-product}{Indexed product and local truth set}
The sort \(s\) of \(M=\prod_{i\in I}M_i\) is
\(\prod_i(M_i)_s\). Constants and functions are interpreted coordinatewise.
For a relation symbol \(R\), including arity zero,
\[
 M\models R(\bar a)\quad\Longleftrightarrow\quad
          (\forall i\in I)\,M_i\models R(\bar a(i)).
\]
For any factor formula \(\psi(\bar x)\) and appropriately sorted product
tuple \(\bar a\), define its local truth set
\[
 [\![\psi(\bar a)]\!]:=
           \{i\in I:M_i\models\psi(\bar a(i))\}\in\mathcal P(I).
\]
The bracket denotes factor satisfaction; it does not assert satisfaction
of \(\psi\) in the product. Free tuples may be empty.
\end{hbdefinition}

\begin{hblemma}{thm:p03-fv-coordinate-terms}{proved=annex:p03-fv-coordinate-terms}{Coordinate evaluation and atomic satisfaction}
Every \(L\)-term evaluates coordinatewise:
\(t^M(\bar a)(i)=t^{M_i}(\bar a(i))\). For every atomic formula
\(\alpha\), including equality, product satisfaction is equivalent to
\([\![\alpha(\bar a)]\!]=I\). This is an atomic fact, not the same
assertion for all first-order formulas.
\end{hblemma}

\begin{hbdefinition}{def:p03-fv-boolean-language}{Boolean conditions and finite partitions}
Use the Boolean-algebra language \(L_B=\{\wedge,\vee,\neg,0,1,=\}\),
interpreted in \(\mathcal P(I)\) by intersection, union, complement,
\(\varnothing,I\). Its first-order quantifiers range over subsets of \(I\).
In particular, this is first-order logic of a Boolean algebra, not merely
a propositional connective combination of truth bits. Write
\(X\le Y\) for the equation \(X\wedge Y=X\).
For a finite set \(F\), define
\[
 \operatorname{Part}_F((Z_f)_{f\in F})
 :=\left(\bigvee_{f\in F}Z_f=1\right)
       \land\bigwedge_{f\ne g}(Z_f\wedge Z_g=0).
\]
Zero cells are permitted. Joins and meets in this display are finite Boolean
terms and finite logical conjunctions respectively. Boolean complement of
a term and logical negation of a Boolean-language formula retain their
different typed roles despite the source's shared \(\neg\) glyph.
\end{hbdefinition}

\section{The effective recursive reduction}

\begin{hbdefinition}{def:p03-fv-reduction}{A reduction sequence and its clauses}
A reduction for \(\varphi(\bar x)\) consists of finitely many
factor formulas \(\psi_1(\bar x),\ldots,\psi_m(\bar x)\) and an
\(L_B\)-formula \(\Theta(X_1,\ldots,X_m)\).
All bound variables are freshened in their own sorts before forming a
clause, and Boolean term substitution is capture-avoiding. In particular,
the newly quantified \(Z_\sigma\) names are fresh for every existing
Boolean binder in \(\Theta\); their repeated schematic names do not
authorize capture across recursive steps. Enumerate bit patterns in a fixed
lexicographic order when forming the finite lists.
An atomic \(\alpha\) uses \(m=1\), \(\psi_1=\alpha\),
\(\Theta(X_1)=(X_1=1)\). Logical falsity uses no factor formulas and
\(\Theta=(0\ne0)\). For logical negation retain the factor list and
negate \(\Theta\) as a formula. For conjunction concatenate the two
factor lists, rename their Boolean variables disjointly, and conjoin their
Boolean conditions. Other Boolean connectives use these clauses; universal
quantification uses \(\forall y\,\chi:=\neg\exists y\,\neg\chi\).

For the decisive existential clause, suppose \(\chi(\bar x,y_s)\)
has reduction \((\psi_1,\ldots,\psi_m;\Theta)\).
For each bit pattern \(\sigma\in\{0,1\}^m\), put
\[
 \delta_\sigma(\bar x,y_s)
  :=\bigwedge_{\sigma_j=1}\psi_j(\bar x,y_s)
                  \land\bigwedge_{\sigma_j=0}\neg\psi_j(\bar x,y_s),
 \qquad
 \eta_\sigma(\bar x):=\exists y_s\,\delta_\sigma(\bar x,y_s).
\]
Empty conjunction is truth. For the new factor list \((\eta_\sigma)_\sigma\),
and Boolean input variables \((Y_\sigma)_\sigma\), define
\[
\begin{split}
 \Theta_\exists((Y_\sigma)_\sigma):=\exists(Z_\sigma)_\sigma\bigg[
 &\operatorname{Part}_{\{0,1\}^m}((Z_\sigma)_\sigma)
       \land\bigwedge_\sigma(Z_\sigma\le Y_\sigma)\\
 &\land\Theta\left(
          \bigvee_{\sigma_1=1}Z_\sigma,\ldots,
          \bigvee_{\sigma_m=1}Z_\sigma\right)\bigg].
\end{split}
\]
Each \(Z_\sigma\) is a Boolean variable, not a product element.
When \(m=0\), there is one empty pattern; the same clause uses its
single partition cell and the existing zero-input Boolean condition.
\end{hbdefinition}

\begin{hbtheorem}{thm:p03-fv-direct-reduction}{proved=annex:p03-fv-direct-reduction}{Feferman--Vaught reduction for every fixed first-order formula}
The recursion terminates and gives a finite reduction sequence for every
fixed \(L\)-formula \(\varphi(\bar x)\), such that for every family
and tuple satisfying the standing assumptions,
\[
 \prod_{i\in I}M_i\models\varphi(\bar a)
 \quad\Longleftrightarrow\quad
 \mathcal P(I)\models\Theta_{\varphi}
 \bigl([\![\psi_1(\bar a)]\!],\ldots,
                    [\![\psi_m(\bar a)]\!]\bigr).
\]
In the existential direction a Boolean partition is converted into an
actual coordinate witness \(b\in\prod_i(M_i)_s\), including the
selected truth pattern at every index. The theorem supplies a syntactic
reduction, not a decision procedure for arbitrary factor theories or a
complexity bound for evaluating the resulting conditions.
\end{hbtheorem}
\begin{hbtheorem}{thm:p03-fv-worked-existential}{proved=annex:p03-fv-worked-existential}{A worked existential with a genuinely global negation}
For binary \(R\) and unary \(P\), set
\[
 U=[\![\exists y(R(x,y)\land\neg P(y))]\!],\qquad
 V=[\![\exists y(R(x,y)\land P(y))]\!].
\]
Then the product satisfies \(\exists y(R(x,y)\land\neg P(y))\)
exactly when \(U\vee V=1\) and \(U\ne0\) in its index Boolean algebra.
The first condition supplies an \(R\)-witness at every index; the second
requires failure of \(P\) at some index, not at every index.
\end{hbtheorem}

\section{Finite-partition Boolean powers}

\begin{hbdefinition}{def:p03-fv-partition-power}{Boolean partitions with base-structure labels}
Let \(A\) be a nonempty \(L\)-structure and \(B\) any Boolean algebra,
possibly trivial. In sort \(s\), a presentation is a finite list
\(((b_r,a_r))_{r=1}^k\), where \(a_r\in A_s\), the \(b_r\) are
pairwise disjoint and \(\bigvee_r b_r=1_B\). Zero coefficients are
allowed; the empty list is permitted precisely when \(0_B=1_B\).
Two presentations \(((b_r,a_r))\), \(((c_t,d_t))\) represent the same
element if
\[
                  a_r\ne d_t\quad\Longrightarrow\quad b_r\wedge c_t=0_B
                    \quad\text{for every }r,t.
\]
Write \(A[B]_{\mathrm{part}}\) for the quotient sorts of these presentations.
This is a finite-partition presentation, to be compared explicitly with the
source's bounded continuous-function power \(A[B]^*\), not silently
identified with a full Cartesian power.

To apply a function symbol to finitely many presented inputs, refine their
partitions by all intersections of input coefficients and label each cell
by the function applied in \(A\) to its input labels. A constant uses the
single cell \(1_B\) with its constant label. For a relation symbol \(R\),
require that the join of refined cells on which the label tuple satisfies
\(R\) in \(A\) is \(1_B\). This includes nullary relations, using the
single empty-input cell.
\end{hbdefinition}

\begin{hbtheorem}{thm:p03-fv-power-structure}{proved=annex:p03-fv-power-structure}{Equivalence, refinement, and formula-valued Boolean elements}
Presentation equality is an equivalence relation and the declared operations
and relations are independent of presentations, so they define an
\(L\)-structure. For every factor formula \(\psi(\bar x)\), define
\([\![\psi(\bar a)]\!]_B\) by refining the parameter partitions and
joining exactly those cells whose label tuple satisfies \(\psi\) in \(A\).
This is a well-defined element of \(B\), even for formulas with quantifiers
in \(A\). It is not defined as truth in the power itself.
\end{hbtheorem}

\begin{hbtheorem}{thm:p03-fv-power-reduction}{proved=annex:p03-fv-power-reduction}{The same reduction with finite Boolean-power witnesses}
For every fixed \(\varphi\), use exactly the reduction sequence above.
Then
\[
 A[B]_{\mathrm{part}}\models\varphi(\bar a)
 \quad\Longleftrightarrow\quad
 B\models\Theta_\varphi
    ([\![\psi_1(\bar a)]\!]_B,\ldots,[\![\psi_m(\bar a)]\!]_B).
\]
The existential witness is a finite partition labeled by finitely many
actual witnesses in \(A\). The construction and proof need only finite
choice, provable in classical \(\mathsf{ZF}\), not an independent selection
at every ultrafilter or index. No Boolean completeness assumption is used.
\end{hbtheorem}

\begin{hbdefinition}{def:p03-fv-represented-power}{A field-of-subsets presentation}
For a Boolean subalgebra \(\mathcal B\subseteq\mathcal P(I)\), the
represented bounded power consists in each sort of finite-range functions
\(a:I\to A_s\) whose fibers \(a^{-1}(d)\) belong to \(\mathcal B\).
Its operations are pointwise and its relations require truth at every index.
Every finite tuple of such functions has a common finite partition into
\(\mathcal B\)-cells on which its values are constant.
\end{hbdefinition}

\begin{hbtheorem}{thm:p03-fv-power-representations}{proved=annex:p03-fv-power-representations}{Finite-range and Stone-space representations}
The represented bounded power is isomorphic to
\(A[\mathcal B]_{\mathrm{part}}\). If \(X\) is a compact Stone space
and \(A_s\) has the discrete topology, a function \(X\to A_s\) is
continuous exactly when it has finite range and clopen fibers. Thus, under
the Boolean Prime Ideal Theorem and the Stone representation of
\cref{thm:p10-stone-representation},
\[
             A[B]_{\mathrm{part}}\cong A[B]^*,
\]
where the source's \(A[B]^*\) consists of continuous functions on \(X_B\)
with pointwise operations. For relational signatures the all-point relation
interpretation is the explicitly proved extension used here. For trivial
\(B\), the Stone space is empty and every power sort has one element.
\end{hbtheorem}

\begin{hbtheorem}{thm:p03-fv-power-elementarity}{proved=annex:p03-fv-power-elementarity}{Exact elementary-preservation consequences}
If \(A\preccurlyeq A'\), the label inclusion gives
\(A[B]_{\mathrm{part}}\preccurlyeq A'[B]_{\mathrm{part}}\).
If \(B\preccurlyeq C\) as Boolean algebras, the coefficient inclusion gives
\(A[B]_{\mathrm{part}}\preccurlyeq A[C]_{\mathrm{part}}\).
Also \(A\equiv A'\) and \(B\equiv B'\) imply
\(A[B]_{\mathrm{part}}\equiv A'[B']_{\mathrm{part}}\).
In particular the finite-range substructure of \(A^I\) is elementary in
the full Cartesian power when \(B=\mathcal P(I)\). The constant diagonal
is not thereby asserted to be elementary.
\end{hbtheorem}

\begin{hbtheorem}{thm:p03-fv-patching-characterization}{proved=annex:p03-fv-patching-characterization}{The source's intrinsic bounded-power criterion}
Let \(X\) be a compact Stone space and \(A\) a nonempty discrete carrier.
A subset \(S\subseteq A^X\) is exactly the set of continuous functions
\(X\to A\) if and only if it contains every constant function,
the equality set of every two of its elements is clopen, and it is closed
under binary patching \(f|_C\cup g|_{X\setminus C}\) for every clopen
\(C\subseteq X\). Compactness changes that binary mixing condition into
the required finite partition construction. This is the carrier-level
criterion of \cite[IV\S5, Theorem~5.6, pp.~143--144]{BurrisSankappanavar2012}.
\end{hbtheorem}

\section{Boundary examples and source distinctions}

\begin{hbtheorem}{thm:p03-fv-boundaries}{proved=annex:p03-fv-boundaries}{What product composition does not permit}
The following boundaries hold.
\begin{enumerate}
\item Product negation is not pointwise universal negation. Two factors can
      make \(R\) true at one coordinate and false at another: the product
      satisfies \(\neg R\), although not every factor does.
\item For a two-element pure equality structure \(A\), the power over
      the two-element Boolean algebra is \(A\), but the power over
      \(\mathcal P(\{0,1\})\) has four elements. The sentence asserting
      three distinct elements separates them, so the constant diagonal
      embedding is not elementary in general.
\item With \(A=\mathbb N\) in the pure equality language and
      \(I=\mathbb N\), finite-range functions form a proper elementary
      substructure of the full power. On the noncompact discrete space
      \(\mathbb N\), the identity function is locally constant but not
      finite-range; compactness cannot be dropped from that equivalence.
\item A Boolean algebra of coefficients is external data in this construction.
      It is not automatically the algebra of all idempotents of a ring power:
      a nonconnected base ring can already contain additional idempotents.
\item Without coordinate witness choice, nonempty product domains alone
      do not justify the existential step of the full-product theorem.
\end{enumerate}
\end{hbtheorem}

\begin{proofboundary}{def:p03-fv-source-boundary}{Direct products, powers, and wider source presentations}
The direct-product Boolean condition, the formal finite-partition power,
and the Stone continuous-function realization have their stated distinct
assumptions. The ring-specific idempotent axioms of the 2023 source are not
asserted for arbitrary \(L\)-structures. Filtered/reduced products,
restricted-support products and enrichments by finite-set predicates require
their own modified constructions; their hypotheses are not borrowed from
the full-product theorem. The complete original Hodges product/Horn,
presentation/word, Feferman--Vaught and Boolean-power sections and their
original examples, exercises and corrections remain separate source duties
\cite[Chapter~9]{Hodges1993}.
\end{proofboundary}
